% Section 10.4, from lectures/L10/appendix/a_flatten_layer.md (Sets 1 to 3),
% lectures/L10/appendix/c_exercises.md (Sets 4 to 6), and the exercise of
% lectures/L10/appendix/d_embedded_constraints.md (Set 7). The solutions to Sets 4 to 6 and to
% Exercise 10.18 are in Appendix A; the solutions to the code are in the repository.

\section{Exercises}
\label{c10:sec:exercises}

\begin{aside}
\textbf{How to work these.} Sets~1 to~3 finish \code{cnn_work}, \file{lectures/L08/cnn_work}: keep
working in the same project, where \file{flatten.hpp} and \file{flatten.cpp} are already in place as
placeholders. Sets~4 to~6 are for self-assessment, reviewing the whole book: Set~4 tests whether you
can recognize a layer type from its algorithm alone, Set~5 tests implementation-level
understanding (shapes, parameter counts, and what has to be remembered for backpropagation), and
Set~6 is a short set of exam-style conceptual questions. Everything you need to answer them is in
the earlier chapters. Set~7 takes the network off the laptop, as \secref{c10:sec:embedded}
describes: one exercise on paper, one in code.

\textbf{How to check your work.} The test suite in \file{lectures/L10/appendix/exercises/test}
points back at \code{cnn_work}; Exercise~\ref{c10:ex:9} runs it, and \code{make ML_DIR=<path>}
points it at a tree of your own elsewhere. The worked solutions to Sets~4 to~6 and to
Exercise~\ref{c10:ex:18} are in Appendix~\ref{sol:ch:solutions}.

The finished reference implementation of the whole pipeline is published in the companion
repository as \file{lectures/L10/cnn_demo}, and this is the chapter it's there for: compare your
\code{cnn_work} against it once your own version runs.
\end{aside}

% ------------------------------------------------------------------------------------------
\exerciseset{Declaring the Flatten Layer}
\label{c10:set:1}\label{c10:app:a}
You'll finish \code{cnn_work} (from Chapters~\ref{c8:ch:conv} and~\ref{c9:ch:maxpool}),
implementing a concrete class \code{Flatten}, replacing the temporary stub class
\code{ml::flatten_layer::Stub} currently used for the flatten layer in \file{factory.cpp}. This is
the same logic you already worked through by hand in \secref{c6:sec:exercises} and implemented as
the standalone struct \code{ml::FlattenLayer} in Chapter~\ref{c7:ch:layers}: here you're adapting
that same logic to satisfy \code{cnn_work}'s polymorphic \code{flatten_layer::Interface}. Once this
is done, every layer in \code{cnn_work} is real, and you'll wire the whole pipeline together for
the first time.

\codeexercise{The \code{Flatten} class, declaration}
\label{c10:ex:1}
In the header file \file{include/ml/flatten_layer/flatten.hpp} (currently a placeholder), add a
class named \code{Flatten} that inherits \code{ml::flatten_layer::Interface}, typeset below
straight from the repository:
\begin{itemize}
  \item Use public inheritance and mark the class \code{final}.
  \item Override every method from the interface, including the destructor (which can be
    \code{= default}).
\end{itemize}

\cppfile{lectures/L08/cnn_work/include/ml/flatten_layer/interface.hpp}

\note Unlike \code{Conv} and \code{MaxPool}, \code{flatten_layer::Interface} declares \textbf{no
\code{optimize()} method at all}: not because it was left out, but because a flatten layer has
nothing to optimize and the interface simply doesn't ask for it. Don't add one.

\codeexercise{Removed constructors and operators}
\label{c10:ex:2}
Delete the default constructor, copy constructor, move constructor, copy assignment operator, and
move assignment operator.

\codeexercise{Private member variables}
\label{c10:ex:3}
Add the following private member variables to \code{Flatten}:
\begin{itemize}
  \item \textbf{\code{myInputGradients}}: the unflattened gradients; what \code{inputGradients()}
    returns (\code{Matrix2d}).
  \item \textbf{\code{myOutput}}: the flattened output; what \code{output()} returns
    (\code{Matrix1d}, not \code{Matrix2d}; note the interface's \code{output()} returns a 1D vector
    here, unlike \code{Conv} and \code{MaxPool}).
\end{itemize}

\codeexercise{Constructor}
\label{c10:ex:4}
Implement the constructor:

\begin{cppcode}
explicit Flatten(std::size_t inputSize) noexcept;
\end{cppcode}

\begin{itemize}
  \item Print an error message and call \code{std::terminate()} if \code{inputSize} is \code{0}.
  \item Size \code{myInputGradients} to \code{inputSize} (a square \code{inputSize × inputSize}
    matrix), and \code{myOutput} to \code{inputSize * inputSize} (see \code{ml::initMatrix()}).
\end{itemize}

\codeexercise{\code{feedforward()}}
\label{c10:ex:5}
Implement \code{bool feedforward(const Matrix2d& input) noexcept}:
\begin{itemize}
  \item Return \code{false} if the size of \code{input} doesn't match the size of
    \code{myInputGradients} (\code{ml::matchDimensions()}), or if \code{input} isn't square
    (\code{ml::isMatrixSquare()}).
  \item Row-major flatten: for every position \code{(i, j)} in \code{input}, write it to
    \code{myOutput[i * inputSize + j]}, where \code{inputSize = myInputGradients.size()}.
  \item Return \code{true}.
\end{itemize}

\codeexercise{\code{backpropagate()}}
\label{c10:ex:6}
Implement \code{bool backpropagate(const Matrix1d& outputGradients) noexcept}:
\begin{itemize}
  \item Return \code{false} if the size of \code{outputGradients} doesn't match the size of
    \code{myOutput}.
  \item The exact inverse of \code{feedforward()}: set
    \code{myInputGradients[i][j] = outputGradients[i * inputSize + j]} for every position
    \code{(i, j)}.
  \item Return \code{true}.
\end{itemize}

\codeexercise{Trivial accessors}
\label{c10:ex:7}
Implement the four accessors from the interface:
\begin{itemize}
  \item \textbf{\code{inputSize()}}: returns \code{myInputGradients.size()}.
  \item \textbf{\code{outputSize()}}: returns \code{myOutput.size()}.
  \item \textbf{\code{inputGradients()}}: returns \code{myInputGradients}.
  \item \textbf{\code{output()}}: returns \code{myOutput}.
\end{itemize}

% ------------------------------------------------------------------------------------------
\exerciseset{Wiring It In}
\label{c10:set:2}

\codeexercise{Wiring the CNN together}
\label{c10:ex:8}
\begin{enumerate}
  \item Make sure \file{source/ml/flatten_layer/flatten.cpp} is listed in \code{cnn_work}'s
    \file{Makefile} under \code{SRC_FILES}. It already is, as it has been since
    Chapter~\ref{c8:ch:conv}: the placeholder compiles to nothing.
  \item In \file{source/ml/factory/factory.cpp}, find \code{Factory::flattenLayer()}: it currently
    returns a \code{flatten_layer::Stub} (marked with a \code{//! @todo} comment). Replace that with
    your new class:
\begin{cppcode}
return std::make_unique<flatten_layer::Flatten>(inputSize);
\end{cppcode}
  \item Run \code{make}. Every layer in \code{cnn_work} (\code{Conv}, \code{MaxPool},
    \code{Flatten}, and \code{Dense}, already done back in Chapter~\ref{c5:ch:dense}) is now real.
    The output of \code{predict()} is finally meaningful, and \code{train()} actually trains a
    working CNN end to end. Compare the predictions against the placeholder output you first saw
    from \code{cnn_work} in Chapter~\ref{c8:ch:conv}.
  \item Compare notes against \file{lectures/L10/cnn_demo}, the finished reference implementation
    sitting right alongside this lecture's material: if your \code{cnn_work} doesn't train the way
    you expect, this is a good place to check your work against.
\end{enumerate}

% ------------------------------------------------------------------------------------------
\exerciseset{Running the Tests}
\label{c10:set:3}

\codeexercise{Running the tests}
\label{c10:ex:9}
The final test suite is available in \file{lectures/L10/appendix/exercises/test}. It's cumulative:
it carries the conv and max pooling layer tests over unchanged and adds unit tests for
\code{Flatten}, plus \textbf{component tests for the network as a whole}.

Build and run it from this lecture's \file{appendix} directory:

\begin{shell}
make -C exercises/test
\end{shell}

\noindent There's nothing to copy: \code{ML_DIR} already points back at the \code{cnn_work}
codebase you started in Chapter~\ref{c8:ch:conv}. All 73 test cases should pass, in about a second.

The component tests are new in this chapter, and they're the reason this suite couldn't exist
earlier: they build a whole network through the real factory and train it, so they only pass once
\code{Conv}, \code{MaxPool} and \code{Flatten} are all real.
\code{LearnsToRecognizeAllFourDigits} trains on the digits 0 to 3 and checks the right output node
wins for each; the closest thing this book has to a final exam for your code.

Two things worth knowing before you start:
\begin{itemize}
  \item The one flatten layer test to read is \code{BackpropagateUndoesFeedforward}. What really
    has to hold is that \code{backpropagate()} undoes exactly what \code{feedforward()} did; if the
    two disagree, gradients land on the wrong pixels and the conv layer quietly learns from noise,
    so the network still runs, it just never gets good. Note though that
    Exercise~\ref{c10:ex:5} mandates \textbf{row-major} order, and \code{FeedforwardIsRowMajor} and
    \code{BackpropagateIsRowMajor} check it, so a self-consistent column-major implementation still
    fails those two.
  \item If \code{LearnsToRecognizeAllFourDigits} fails but every unit test passes, the layers are
    individually right and something about how they're wired together is wrong. Start with
    \code{BackpropagationReachesTheConvLayerThroughEveryLayer}, which chains the four layers by hand
    and checks the gradients survive the trip back.
\end{itemize}
The test suite's README, \file{lectures/L10/appendix/exercises/test/README.md}, has more
information, including how the convergence thresholds were measured.

% ------------------------------------------------------------------------------------------
\exerciseset{Recognize the Layer}
\label{c10:set:4}\label{c10:app:c}

\reflexercise{Six feedforward algorithms}
\label{c10:ex:10}
Below are six feedforward algorithms, written as pseudocode. For each one, decide which layer type
it describes: a linear regression layer, a dense layer, a conv layer, a max pooling layer, or a
flatten layer.

\task{A}
\begin{pycode}
output = k * input + m
\end{pycode}

\task{B}
\begin{pycode}
for n in nodes:
    sum = bias[n]
    for i in inputs:
        sum += weight[n][i] * input[i]
    output[n] = activation(sum)
\end{pycode}

\task{C}
\begin{pycode}
for row in range(outputHeight):
    for col in range(outputWidth):
        sum = bias
        for kr in range(kernelSize):
            for kc in range(kernelSize):
                sum += kernel[kr][kc] * input[row + kr][col + kc]
        output[row][col] = activation(sum)
\end{pycode}

\task{D}
\begin{pycode}
for row in range(outputHeight):
    for col in range(outputWidth):
        window = input[row * poolSize : row * poolSize + poolSize,
                        col * poolSize : col * poolSize + poolSize]
        output[row][col] = max(window)
\end{pycode}

\task{E}
\begin{pycode}
output = reshape(input, (1, height * width))
\end{pycode}

\task{F}
\begin{pycode}
sum = bias[0]
for i in inputs:
    sum += weight[0][i] * input[i]
output[0] = sum   // no activation function applied
\end{pycode}
\seesolution{c10:ex:10}

% ------------------------------------------------------------------------------------------
\exerciseset{Implementation Questions}
\label{c10:set:5}
Consider the following network, processing a single-channel \textbf{8\,×\,8} input image:

\begin{console}
input (8×8)
  → conv layer (one 3×3 kernel, stride 1, zero-padded)
  → max pooling layer (2×2, stride 2)
  → flatten layer
  → dense layer (10 nodes)
\end{console}

\calcexercise{Tracing shapes}
\label{c10:ex:11}
What is the shape of the data after each layer (conv, max pooling, flatten, dense)?
\seesolution{c10:ex:11}

\calcexercise{Counting trainable parameters}
\label{c10:ex:12}
How many trainable parameters does each layer have? What is the total for the whole network?
\seesolution{c10:ex:12}

\reflexercise{What max pooling has to remember}
\label{c10:ex:13}
The max pooling layer has no trainable parameters, yet it still needs to store something during
\code{feedforward()} so that \code{backpropagate()} works correctly. What does it need to remember,
and why?
\seesolution{c10:ex:13}

\calcexercise{Conv vs.\ dense, parameter count}
\label{c10:ex:14}
A dense layer connecting the flattened 8\,×\,8 image (64 values) directly to a single output node
would need 64 weights + 1 bias = 65 trainable parameters just to combine the whole image once. The
conv layer in the network above combines local pixel neighborhoods using only how many trainable
parameters? Why is the difference so large?
\seesolution{c10:ex:14}

\reflexercise{Flatten in reverse}
\label{c10:ex:15}
During \code{feedforward()}, a flatten layer reshapes a \textbf{3\,×\,4} matrix into a
\textbf{1\,×\,12} vector. During \code{backpropagate()}, it receives a \textbf{1\,×\,12} gradient
vector from the next layer. What shape does it reshape that gradient vector back into, and why?
\seesolution{c10:ex:15}

% ------------------------------------------------------------------------------------------
\exerciseset{Exam-Style Questions}
\label{c10:set:6}

\reflexercise{True or false?}
\label{c10:ex:16}
For each statement, decide if it's true or false, and briefly justify your answer.
\begin{enumerate}
  \item A max pooling layer has trainable parameters that are updated during \code{optimize()}.
  \item The same kernel (the same weights) is reused at every spatial position in a conv layer.
  \item A flatten layer's \code{backpropagate()} performs the exact same reshape operation as its
    \code{feedforward()}.
  \item A dense layer with exactly one node, no activation function, and no hidden layer around it
    behaves identically to a linear regression layer.
  \item Adding more filters (kernels) to a conv layer increases its trainable parameter count, but
    doesn't by itself change the width or height of its output.
\end{enumerate}
\seesolution{c10:ex:16}

\reflexercise{Short answer}
\label{c10:ex:17}
\begin{enumerate}[start=6]
  \item In one or two sentences, explain why swapping a dense layer for a conv layer generally
    means far fewer trainable parameters for the same size of input image.
  \item A student claims: \emph{``Max pooling layers slow down training because they have so many
    weights to update.''} Is this correct? Why or why not?
  \item Most layer types in this book implement the same three methods: \code{feedforward()},
    \code{backpropagate()}, and \code{optimize()}. What does \code{optimize()} do for a max pooling
    layer, which has no trainable parameters? And why does the flatten layer's interface not
    declare \code{optimize()} at all, rather than declaring it and leaving it empty?
\end{enumerate}
\seesolution{c10:ex:17}

% ------------------------------------------------------------------------------------------
\exerciseset{Off the Laptop}
\label{c10:set:7}
Two of the constraints from \secref{c10:sec:embedded}, worked through by hand: memory footprint on
paper, and fixed-point arithmetic in code.

\calcexercise{Memory footprint}
\label{c10:ex:18}
Using the layer sizes from the network you built in Chapters~\ref{c4:ch:network}
and~\ref{c5:ch:dense} (or any other layer configuration of your choice):
\begin{enumerate}
  \item Compute the total number of trainable values (weights and biases) across all layers.
  \item Compute the memory footprint at \code{double} (8 bytes), \code{float} (4 bytes), and
    \code{int8} (1 byte) per value.
  \item Compare your numbers against the example in \secref{c10:sec:memory} of a layer sized for
    28\,×\,28 pixel images (785\,KB, 392.5\,KB and 98.125\,KB). Which of the three precisions would
    plausibly fit into a microcontroller with 64\,KB of RAM? With 512\,KB?
\end{enumerate}
\seesolution{c10:ex:18}

\codeexercise{Fixed-point conversion}
\label{c10:ex:19}
Take the \code{predict()} method from Chapter~\ref{c1:ch:linreg}'s \code{ml::lin_reg::Fixed} class
($y = k \cdot x + m$):
\begin{enumerate}
  \item Add \code{toFixed()}, \code{toDouble()}, and \code{fixedMultiply()} as shown in
    \secref{c10:sec:fixedpoint} (an addition helper isn't needed: plain integer \code{+} already
    works for Q16.16 values).
  \item Implement a free function
    \code{Fixed16_16 predictFixed(Fixed16_16 weight, Fixed16_16 bias, Fixed16_16 input)} that
    computes the same prediction as \code{predict()}, but entirely in Q16.16 fixed-point using
    \code{fixedMultiply()} for the multiplication.
  \item In \code{main()}, pick a few \code{(weight, bias, input)} triples, convert them to
    fixed-point, run both \code{predict()} (floating-point) and \code{predictFixed()}
    (fixed-point), convert the fixed-point result back with \code{toDouble()}, and print both side
    by side.
  \item Confirm the two results agree to within roughly 0.0001 for each test case. If they don't,
    double-check the \code{>>} shift in \code{fixedMultiply()}: it's the step most likely to be off
    by a factor of \code{kOne}.
\end{enumerate}
